\subsection{\mscre -- Physik des Kontinuums}\label{kap:kontmech_matlabskripts}

Die \mscre befinden sich im MATLAB-Ordner \texttt{MatlabDateien/kontmech},
die Hilfsfunktionen im dortigen Unterordner \texttt{Include\-folder} und Daten im Unterordner \texttt{Data}.
Es wird empfohlen, entweder die Skripte im GitHub Verzeichnis \githubSN zu benutzen oder sich die komplette Verzeichnisstruktur von \MKc~ runter zu laden und diese und auch das Arbeitsbuch im Hauptverzeichnis Fingeruebungen zu installieren. \\

Wir wollen noch einmal explizit darauf hinweisen, dass die \mscre prim�re der Vermittlung der physikalischen Modelle dienen und keinen Anspruch auf besonders elegante oder effiziente Programmierung erheben. Unser Buch ist ja kein Lehrbuch f�r \matlab-Programmierung. Uns ging es bewusst darum, dass die Skripte gut lesbar sind und das physikalische Modell bzw den physikalischen Gedankengang m�glichst klar abbilden.
Die Leser sind aufgerufen, die Programme nach Belieben zu verbessern und kompakter zu gestalten und diese auch ins GitHub Verzeichnis \githubSN zu stellen.


%\begin{prob}\label{mls:kontmech_hilfsfunktionen} %MLSkript1 
%\textbf{Hilfsfunktionen} \\
%\begin{longtable}{p{0.33\textwidth}p{0.67\textwidth}}
%\hline\noalign{\smallskip}
%Skript & Beschreibung \\
%\noalign{\smallskip}\svhline\noalign{\smallskip}
%Mathematische Funktionen \\
%\noalign{\smallskip}\svhline\noalign{\smallskip}
%\mlkontmechref{Beispiel.m} & Ein Beispieltext.\\
%\noalign{\smallskip}\svhline\noalign{\smallskip}
%\end{longtable}
%\end{prob}

\subsubsection {Skripte zu Beispielen} \label{mls:kontmech_kapitel} 
\begin{longtable}{p{0.33\textwidth}p{0.67\textwidth}}
\hline\noalign{\smallskip}
Skript & Beschreibung \\
\noalign{\smallskip}\svhline\noalign{\smallskip}
\mlkontmechref{SchwingendeSaite01.m} & L�sen der Differentialgleichung einer
schwingenden Saite f�r eine stehende Welle.\\
\noalign{\smallskip}\svhline\noalign{\smallskip}
\end{longtable}

\subsubsection {Skripte zu �bungen} \label{mls:kontmech_uebungen} 
\begin{longtable}{p{0.33\textwidth}p{0.67\textwidth}}
\hline\noalign{\smallskip}
Skript & Beschreibung \\
\noalign{\smallskip}\svhline\noalign{\smallskip}
\mlkontmechref{AchterbahnKont.m}        & L�sen der Differentialgleichung eines kontinuierlichen, elastischen Achterbahnzuges auf einer Gau�kurven-Strecke.\\
\mlkontmechref{AchterbahnKontBeschl.m}  & Normalbeschleunigung auf einen Fahrgast im kontinuierlichen, elastischen Achterbahnzug.\\
\mlkontmechref{Chladni01.m}             & L�sen der Eigenschwingungen einer kreisf�rmigen, elastischen Fl�che.\\
\mlkontmechref{Chladni02.m}             & L�sen der Eigenschwingungen einer quadratischen, elastischen Fl�che.\\
\mlkontmechref{Chladni03.m}             & L�sen der Eigenschwingungen einer rechteckigen, elastischen Fl�che.\\
\mlkontmechref{Divergenz.m}             & Programm berechnet Vektorfunktionen und ihre Divergenz.\\
\mlkontmechref{FassFormelGauss.m}       & Programm berechnet Volumen von Rotationsk�rpern �ber Gau�schen Integralsatz.\\
\mlkontmechref{GespannterBogen.m}       & L�sen der Eigenschwingungen eines Bogens, der durch eine Saite gespannt ist.\\
\mlkontmechref{GradientFunktion1.m}     & Programm berechnet Funktionen und ihre Gradienten.\\
\mlkontmechref{GradientFunktion2.m}     & Programm berechnet Funktionen und ihre Gradienten (Yukawa-Potential).\\
\mlkontmechref{GradientFunktion3.m}     & Programm berechnet Funktionen und ihre Gradienten (Kepler-Potential mit Drehimpuls).\\
\mlkontmechref{KdVTsunami.m}            & Berechnung eines KdV-Solitons, das auf eine K�ste zul�uft, an der die Wassertiefe abnimmt.\\
\mlkontmechref{KdV01.m}                 & L�sen der Korteweg-de-Vries-Gleichung f�r ein einzelnes Soliton.\\
\mlkontmechref{KdV02.m}                 & Vergleich mehrerer Solitonen der KdV-Gleichung mit unterschiedlicher Geschwindigkeit.\\
\mlkontmechref{KdV03.m}                 & Summen aus einzelnen Solitonen in der KdV-Gleichung.\\
\mlkontmechref{KdV04.m}                 & Interaktion von Solitonen in der KdV-Gleichung.\\
\mlkontmechref{Kreisrohr.m}             & Berechnen des Geschwindigkeitsprofils einer station�ren Str�mung in einem kreisrunden Rohr.\\
\mlkontmechref{NLSGSolitonen.m}         & L�sen der nichtlinearen Schr�dinger-Gleichung f�r ein Soliton.\\
\mlkontmechref{Rechteckrohr.m}          & Berechnen des Geschwindigkeitsprofils einer station�ren Str�mung in einem rechteckigen Rohr.\\
\mlkontmechref{Rotation.m}              & Programm berechnet Vektorfunktionen und ihre Rotation.\\
\mlkontmechref{SchwingendeSaite02.m}    & L�sen der Differentialgleichung einer eingespannten Saite f�r eine Gau�sche St�rung.\\
\mlkontmechref{SchwingendeSaite03.m}    & L�sen der Differentialgleichung einer eingespannten Saite f�r eine einseitig angeregte Welle.\\
%\noalign{\smallskip}\svhline\noalign{\smallskip}
\mlkontmechref{SinusGordonSolitonen.m}  & L�sen der Sinus-Gordon-Gleichung f�r ein Soliton.\\
\mlkontmechref{Tank.m}  & Fluss durch einen quadratischen Tank mit zwei �ffnungen.\\
\mlkontmechref{Torricelli.m}  & Fluss einer Fl�ssigkeit aus einem Tank nach dem
Gesetz von Torricelli.\\
\end{longtable}



